\documentclass[aspectratio=169,10pt]{beamer}
\usepackage{default_latex_beamer_template}
\usepackage{amsmath}
\hypersetup{pdftitle={NYgrid 2019 interface-flow verification and compact-model calibration diagnostics},pdfauthor={NPCC / New York modeling project},colorlinks=true,urlcolor=blue,linkcolor=muted}
\pdfstringdefDisableCommands{\def\translate#1{#1}}
\setbeamertemplate{itemize item}{\color{teal}\small\textbullet}
\setlength{\tabcolsep}{7pt}
\renewcommand{\arraystretch}{1.16}
\renewcommand{\source}[1]{\par\vfill{\color{muted}\fontsize{6.5}{7.5}\selectfont Source: #1\par}}
\newcommand{\lead}[1]{{\large\color{navy}\bfseries #1}\par\medskip}
\newcommand{\minor}[1]{{\small\color{muted}#1}\par}
\newcommand{\tightitems}{\setlength{\itemsep}{0.65em}}
\newcommand{\plot}[2]{\begin{center}\includegraphics[height=#2,keepaspectratio]{figures/#1}\end{center}}
\title[NYgrid 2019 | Interface-flow project update]{NYgrid interface-flow verification\\and compact-model calibration diagnostics}
\subtitle{Historical 2019 reconstruction and preliminary DLR research}
\author{NPCC / New York modeling project}
\date{September 8, 2026}

\begin{document}

% 1
\defaulttitleframe{Project update: released NYgrid study, matched-input experiments, selected snapshots, and remaining validation}

% 2
\begin{frame}{Annual accuracy improves, with a clear interface tradeoff}
\lead{The compact model's annual pooled error falls by 5.07 percentage points.}
\begin{center}
\begin{tabular}{lrrr}
\toprule
Annual WAPE & NYgrid reference & Compact 51 & Compact 51 + Marcy \\
\midrule
All seven interfaces & 11.22\% & 17.38\% & \textbf{12.32\%} \\
Central East & 9.63\% & 49.23\% & \textcolor{teal}{\textbf{7.09\%}} \\
UPNY--ConEd & 9.91\% & 10.55\% & \textcolor{red}{\textbf{18.59\%}} \\
\bottomrule
\end{tabular}
\end{center}
\par\medskip
\begin{itemize}\tightitems
\item All three series use the same 8,760 hours, 2019 source inputs, and corrected paper-role interface definitions.
\item A selected compact hour reaches \textbf{3.28\% pooled error}, with every interface below \textbf{4.86\%}.
\item This experiment diagnoses representation and allocation choices; shared-parameter fitting and untouched predictive evaluation remain next steps.
\end{itemize}
\source{[3,4]; frozen annual/monthly metrics and best-snapshot records, project commit c472e25. WAPE uses actual flow as denominator.}
\end{frame}

% 3
\begin{frame}{The low-error results belong to a historical 51-bus variant}
\lead{The existing 71-bus year-end 2025 model was unchanged.}
\small
\renewcommand{\arraystretch}{1.08}
\begin{tabularx}{\textwidth}{lrr>{\raggedright\arraybackslash}X}
\toprule
Network role & Buses & AC branches\textsuperscript{*} & Purpose in this update \\
\midrule
Released NYgrid source & 57 & 94 & 2019 DC study, including external equivalents. \\
NY-only reference control & 46 & 67 & Exact snapshot replay with fixed NY-side boundary injections. \\
Historical compact & 51 & 92 / 87 & Pre-2025 parent with matched 2019 injections. \\
Compact + Marcy hypothesis & 51 & 94 / 89 & Two added source-informed branches. \\
Year-end 2025 variant & 71 & 135 / 117 & Existing expansion; unchanged here. \\
\bottomrule
\end{tabularx}
\par\medskip
\minor{\textsuperscript{*}Stored / active where shown. The source also has four separate DC-link records.}
\par\medskip
The separate AC checkpoint uses the \textbf{51-bus network without Marcy}, with two restored nuclear proxies and 37 generator records.
\source{[3,4]; reference export contract, compact export description, and model ledger. Historical compact is a research reconstruction, not an exact 2019 as-built system.}
\end{frame}

% 4
\begin{frame}{NYgrid combines network reduction with reconstructed operation}
\begin{columns}[T,onlytextwidth]
\column{0.45\textwidth}
\lead{Network representation}
\begin{itemize}\tightitems
\item Start from the 140-bus NPCC benchmark.
\item Modified Ward reduction preserves 46 New York buses, nine external boundary buses, and two auxiliaries.
\item Add the paper's selected network modifications and controlled transfers.
\end{itemize}
\column{0.51\textwidth}
\lead{2019 operating inputs}
\begin{itemize}\tightitems
\item Thermal profiles and daily nuclear reconstruction.
\item Hourly hydro split: 80\% Niagara / St. Lawrence, 20\% smaller hydro; monthly St. Lawrence output.
\item NYISO zonal demand and external schedules.
\item Released generation-residual policy assigns uncovered generation to Zone J.
\end{itemize}
\end{columns}
\par\bigskip
\textbf{Then solve DC power flow and compare interface flows with NYISO.}\par
\minor{The paper was published in 2023; the operating study uses 2019 data. Reconstructed generation is not independent unit telemetry, and the flow replay does not enforce operating bounds.}
\source{[1,2,3]; Liu et al., IEEE Transactions on Power Systems 38(4), 2023; released updateOpCond, allocateLoad, and modifyMPC code.}
\end{frame}

% 5
\begin{frame}{``Percent error'' changes meaning with its denominator}
\small
\begin{tabularx}{\textwidth}{>{\raggedright\arraybackslash}p{3.1cm}p{5.3cm}>{\raggedright\arraybackslash}X}
\toprule
Metric & Definition & Interpretation \\
\midrule
Paper Eq. 11 & $100\,(F_{\rm obs}-F_{\rm model})/L^{+}$ & Signed mismatch divided by the hourly positive limit. \\[1ex]
Selected-hour error & $100\,|F_{\rm model}-F_{\rm obs}|/|F_{\rm obs}|$ & Proportional error against that interface's actual flow. \\[1ex]
Annual interface WAPE & $100\,\dfrac{\sum_t |F_{{\rm model},t}-F_{{\rm obs},t}|}{\sum_t |F_{{\rm obs},t}|}$ & Annual absolute mismatch divided by annual absolute observed flow. \\
\bottomrule
\end{tabularx}
\par\medskip
\begin{itemize}\tightitems
\item The paper's ``within 10\%'' statement is \textbf{limit-normalized}; it does not mean 10\% of actual flow.
\item West Central uses a positive-limit value of 9999 throughout the released data; this is not a verified physical rating.
\item Pooled WAPE sums across seven overlapping interfaces. It measures comparison accuracy, not statewide energy imbalance.
\end{itemize}
\source{[1,2,3]; paper Eq. 11, released flow4Plot code, and metric-definition audit. Near-zero observed flows can inflate selected-hour proportional errors.}
\end{frame}

% 6
\begin{frame}{Repairing the measurement formulas precedes calibration}
\lead{Five stale plotting formulas selected the wrong branch rows.}
\begin{center}
\begin{tikzpicture}[>=Latex, every node/.style={font=\small,align=center}]
\node[draw=navy,rounded corners,minimum width=3.0cm,minimum height=1.1cm] (pf) at (0,0) {Same 8,760\\solved electrical states};
\node[anchor=west] (old) at (5,0.7) {Literal plotting rows\\\textcolor{red}{\textbf{60.66\% pooled WAPE}}};
\node[anchor=west] (new) at (5,-0.8) {Corrected released interface map\\\textcolor{teal}{\textbf{11.22\% pooled WAPE}}};
\draw[->,thick,muted] (pf.east) -- ++(1.0,0) |- (old.west);
\draw[->,thick,teal] (pf.east) -- ++(1.0,0) |- (new.west);
\end{tikzpicture}
\end{center}
\small
Affected: Dysinger East, West Central, Moses South, UPNY--ConEd, and Dunwoodie/Sprain. Central East and Total East were already correct.
\par\medskip
\textbf{UPNY example:} the old expression included bus 71--72, a C--C branch. The corrected map uses the directed G--H crossings 73--74 and 77--74.
\par\medskip
\minor{Endpoint and role checks support the repair. No injection, impedance, dispatch, or observed target was fitted in this step.}
\source{[2,3]; operator audit formulas and terms; annual reproduction metrics. This is an audited replay of released data/code, with a documented measurement correction.}
\end{frame}

% 7
\begin{frame}{The reference is strongest on major eastern interfaces}
\plot{01_nygrid_annual_interface_wape.pdf}{5.45cm}
\minor{Annual pooled WAPE: \textbf{11.22\%}. West Central is the weakest interface: 46.81\% WAPE and 1,006 opposite-direction hours.}
\source{[3]; corrected released interface map, all 8,760 prepared hours in 2019. Total seven-interface opposite-direction count: 1,092 / 61,320.}
\end{frame}

% 8
\begin{frame}{NYgrid has a strong selected hour across all seven interfaces}
\lead{January 8, 2019 at 15:00 New York local time}
\small
\begin{center}
\begin{tabular}{lrrrr}
\toprule
Interface & Observed MW & Model MW & Absolute error MW & Error \\
\midrule
Dysinger East & 1,471.7 & 1,467.3 & 4.3 & 0.30\% \\
West Central & 844.6 & 875.3 & 30.7 & 3.63\% \\
Moses South & 1,996.8 & 2,078.0 & 81.2 & 4.07\% \\
Central East & 2,725.4 & 2,772.5 & 47.1 & 1.73\% \\
Total East & 4,738.2 & 4,856.4 & 118.1 & 2.49\% \\
UPNY--ConEd & 2,106.4 & 2,153.2 & 46.7 & 2.22\% \\
Dunwoodie/Sprain & 2,422.6 & 2,419.5 & 3.1 & 0.13\% \\
\bottomrule
\end{tabular}
\end{center}
\textbf{Pooled: 2.03\%. Worst interface: 4.07\%. No opposite directions.}
\par\medskip
\minor{Best under both pooled-error and minimax ranking, selected retrospectively from the examined year. All seven channels have 12 distinct samples and no imputation or DST ambiguity.}
\source{[3]; output/nygrid\_2019/best\_snapshots/BEST\_SNAPSHOTS.md and selected per-interface records. These are actual-scale MW. Annual error remains 11.22\%.}
\end{frame}

% 9
\begin{frame}{The matched experiment separates network and input effects}
\begin{center}
\begin{tabular}{lcc}
\toprule
 & Released paper inputs & Current pipeline's 2019 inputs \\
\midrule
46-bus NYgrid control & A & B \\
51-bus compact & C & D \\
51-bus compact + Marcy & C + Marcy & D + Marcy \\
\bottomrule
\end{tabular}
\end{center}
\begin{itemize}\tightitems
\item \textbf{Annual comparison:} A, C, and C + Marcy across all 8,760 hours. The six-arm input/network matrix uses five common development hours.
\item \textbf{One power scale:} $\gamma=0.3586622254$ multiplies load, generation, boundary injection, and flow targets. Reactances are unchanged; WAPE is invariant.
\item \textbf{Consistent comparison:} paper-role interface cuts and bus 74 as the DC balancing reference. Native input generation, solved generation, and slack adjustment are recorded separately.
\item \textbf{Exact control:} export NY-side AC crossings, four controlled links, and HQ once. The 46-bus control reproduces source internal flows to approximately $10^{-11}$ MW.
\end{itemize}
\source{[4]; reference export, network/input matrix, and independent validation. Boundary values are model-derived terminal injections; exactness applies to each fixed snapshot.}
\end{frame}

% 10
\begin{frame}{Seven bus-role differences affect allocation and metering}
\begin{columns}[T,onlytextwidth]
\column{0.54\textwidth}
\small
\begin{tabular}{rlcc}
\toprule
Bus & Name & Paper role & Compact role \\
\midrule
38 & Gilboa & E & F \\
46 & Colton & E & D \\
47 & Moses W & E & D \\
62 & Stolle Road & B & A \\
69 & Binghamton & E & C \\
77 & Buchanan & G & H \\
79 & RAVA-3 & K & J \\
\bottomrule
\end{tabular}
\column{0.43\textwidth}
\lead{Use a common role system to interpret the comparison.}
\small
All 46 original NY buses remain in compact51. Source generation at buses 45, 49, 62, 74, and 77 can be preserved as DC bus injections despite absent compact generator records.
\par\medskip
Restored Indian Point proxies land at source bus 74 in the separate AC fleet experiment.
\end{columns}
\par\bigskip
\textbf{Sensitivity:} applying current physical-zone cuts to the same annual states gives pooled WAPE \textbf{19.07\% $\rightarrow$ 19.67\%} after Marcy. Operator meaning materially changes the score.
\source{[4]; bus-role and operator crosswalks; annual metrics. Roles are equivalent-model assignments, not a claim that one set of labels is geographic truth.}
\end{frame}

% 11
\begin{frame}{The Marcy experiment tests two specific network branches}
\begin{center}
\begin{tikzpicture}[>=Latex, every node/.style={font=\small,align=center},bus/.style={circle,draw=navy,thick,minimum size=0.8cm,fill=white}]
\node[bus] (a) at (0,0.3) {43};
\node[bus] (b) at (4.0,0.3) {38};
\node[bus] (c) at (8.0,0.3) {77};
\node[bus] (d) at (2.2,-1.7) {69};
\node[bus] (e) at (5.8,-1.7) {39};
\draw[teal,line width=1.7pt] (a) -- node[above=3pt]{Added: $X=0.0427$ pu} (b);
\draw[teal,line width=1.7pt] (b) -- node[above=3pt]{Added: $X=0.0147$ pu} (c);
\draw[muted,thick] (b) -- (d);
\draw[muted,thick] (b) -- (e);
\node[anchor=east,text=muted] at (2.1,-0.8) {Retained\\$X=0.0485$ pu};
\node[anchor=west,text=muted] at (5.9,-0.8) {Retained S4\\Gilboa--Leeds};
\end{tikzpicture}
\end{center}
\begin{itemize}\tightitems
\item Two \textbf{AC-network branches used in the DC approximation}, with reactances from the paper's model on a 100 MVA base. Bus count stays at 51.
\item The inherited compact backbone and its parallel paths remain. No identical active endpoint pair is duplicated; broader functional overlap is unresolved.
\item $R=B=0$ and RATE\_A = 0 are numerical placeholders. Physical parameters, ratings, AC controls, and DLR suitability remain unqualified.
\end{itemize}
\source{[2,4]; modifyMPC, Marcy hypothesis ledger, and network-pair audit. Schematic shows modeled connectivity only; positions are not geographic.}
\end{frame}

% 12
\begin{frame}{Central East improves, but UPNY--ConEd gets worse}
\plot{02_annual_model_comparison.pdf}{5.55cm}
\minor{Compact pooled WAPE: \textbf{17.38\% $\rightarrow$ 12.32\%}. West Central remains at 46.81\%; the variant does not meet a uniform per-interface accuracy target.}
\source{[4]; annual metrics, all 8,760 hours with identical paper inputs and paper-role cuts. Marcy branch parameters are source-informed, not fitted to these observations.}
\end{frame}

% 13
\begin{frame}{The pooled improvement appears in every month}
\plot{03_monthly_pooled_wape.pdf}{5.55cm}
\minor{Monthly improvement ranges from approximately \textbf{4.06 to 5.87 percentage points}. These are all examined months, not untouched validation blocks.}
\source{[4]; annual\_monthly\_metrics.csv. Common 2019 inputs and paper-role operators across all three network series.}
\end{frame}

% 14
\begin{frame}{Total East reveals regional balance rather than path sharing}
\lead{For a complete cut in the fixed-injection, lossless DC problem:}
\[
\underbrace{F_{\rm Total\ East}}_{\text{western export}}
=\sum_{i\in\mathcal{W}}\left(P_{G,i}+P_{{\rm boundary},i}-P_{{\rm net\ load},i}\right)
\]
\begin{columns}[T,onlytextwidth]
\column{0.47\textwidth}
\textbf{What Marcy changes}\par\medskip
Reactances redistribute flow among parallel paths. Central East WAPE falls from 49.23\% to 7.09\%.
\column{0.47\textwidth}
\textbf{What remains fixed}\par\medskip
The complete regional net export is set by injections. Total East WAPE stays at 7.235\% in the matched experiment.
\end{columns}
\par\bigskip
\textbf{Calibration implication:} regional balance errors require checking load, generation, boundary allocation, and interface membership. Impedance changes alone cannot remove a complete-cut mismatch with fixed injections.
\source{[4]; Total East regional ledger and annual metrics. This identity uses the study's complete western cut; AC losses and redispatch require additional terms.}
\end{frame}

% 15
\begin{frame}{Allocation tests identify candidates, not a universal policy}
\lead{Five common development hours, compact + Marcy network}
\small
\begin{tabularx}{\textwidth}{lc>{\raggedright\arraybackslash}X}
\toprule
Change from released policy & Hours improved & Interpretation \\
\midrule
Statewide thermal residual & 4 / 5 & The released Zone-J-only policy is not consistently superior. \\
Current hydro spatial rule & 2 / 5 & Hydro location needs better evidence; neither tested rule wins throughout. \\
Current load roles + weights & 0 / 5 & A combined spatial reassignment can strongly worsen flows despite retained demand totals. \\
Current boundary preparation & 4 / 5 & Boundary allocation and time aggregation are material sources of mismatch. \\
\bottomrule
\end{tabularx}
\par\bigskip
\textbf{Keep the interpretation controlled:} the load test changes roles and weights together; the boundary test also changes preparation methodology and can change net imports/slack.
\par\medskip
\minor{No delivered model switches to the best policy hour by hour. These are fixed-policy comparisons to guide the next shared-parameter calibration.}
\source{[4]; allocation\_ablations.csv and network/input matrix. Five hours are development evidence; no new target-fitting or held-out test was performed.}
\end{frame}

% 16
\begin{frame}{Three compact hours keep every interface error near 5\%}
\small
\begin{center}
\begin{tabular}{lrrr}
\toprule
2019 local hour & Jul 3, 04:00 & Jul 22, 04:00 & Sep 8, 05:00 \\
\midrule
Dysinger East & 1.46\% & 0.44\% & 2.57\% \\
West Central & 1.48\% & 5.31\% & 3.14\% \\
Moses South & 4.85\% & 2.26\% & 0.96\% \\
Central East & 4.22\% & 3.39\% & 4.32\% \\
Total East & 3.02\% & 1.42\% & 0.40\% \\
UPNY--ConEd & 3.75\% & 1.62\% & 5.59\% \\
Dunwoodie/Sprain & 2.45\% & 5.09\% & 0.01\% \\
\midrule
\textbf{Pooled error} & \textbf{3.28\%} & \textbf{2.63\%} & \textbf{1.95\%} \\
\textbf{Worst interface} & \textbf{4.85\%} & \textbf{5.31\%} & \textbf{5.59\%} \\
\bottomrule
\end{tabular}
\end{center}
\minor{Top three \textbf{minimax} hours: rank each hour by its largest individual proportional error. No opposite directions; no imputation or DST ambiguity; 12 distinct samples per channel.}
\source{[4]; compact + Marcy, paper inputs and paper-role cuts. Retrospective development examples selected from the examined year; no per-hour target fit.}
\end{frame}

% 17
\begin{frame}{Selected-hour flows closely track the observed magnitudes}
\plot{04_selected_snapshot_flows.pdf}{5.55cm}
\minor{Two independently selected dates, \textbf{not a same-hour network comparison}. Compact flows are divided by $\gamma$ for an actual-scale GW display; the operating model remains at benchmark scale.}
\source{[3,4]; source Jan 8 best hour and compact + Marcy Jul 3 best minimax hour. Actual-flow proportional errors; retrospective selection.}
\end{frame}

% 18
\begin{frame}{A separate AC checkpoint establishes three feasible examples}
\lead{Unchanged compact51 network, restored 2019 fleet, current-pipeline inputs}
\small
\begin{tabular}{lrr}
\toprule
Three-hour pooled WAPE & DC & Bounded AC + fresh PF \\
\midrule
Current physical-zone operators & 19.06\% & 17.54\% \\
Paper-role operators & 23.43\% & 22.47\% \\
\bottomrule
\end{tabular}
\par\medskip
\begin{itemize}\tightitems
\item January 21, April 21, and July 20 pass the bounded AC solve and a fresh power-flow replay. No interface-flow objective or relaxed operating limits.
\item These cases use \textbf{no Marcy additions} and AC reference bus 42. They do not establish AC feasibility of the best Marcy DC hours.
\item Feasible controls differ materially from the raw prior. April required an IPOPT fallback; losses rose from \textbf{98.99 to 206.54 benchmark MW} under the control adjustment.
\end{itemize}
\minor{The result establishes bounded feasible operating examples, not reconstructed historical voltage/reactive controls. Some individual interface errors worsen.}
\source{[4]; AC feasibility/control summary and paired scoring. Three seasonal hours, 21 interface comparisons per operator set; AC uses terminal PF/PT consistently.}
\end{frame}

% 19
\begin{frame}{Next, turn the diagnostics into a validated calibration}
\begin{enumerate}\tightitems
\item \textbf{Freeze the model and measurement contract.}\par
\small Decide intended allocation roles and interface cuts; keep original bus IDs, boundary signs, and power scaling traceable.\normalsize
\item \textbf{Fit a small, physically bounded set of shared parameters.}\par
\small Prioritize UPNY sharing, West Central direction/error, hydro allocation, and uncovered generation. Avoid per-hour corrections that absorb every mismatch.\normalsize
\item \textbf{Build a credible AC realization of the candidate topology.}\par
\small Resolve Marcy representation, finite line parameters/ratings, parallel-path overlap, generator capability, voltage and reactive controls.\normalsize
\item \textbf{Evaluate on data not used to choose or fit the model.}\par
\small Predeclare metrics and holdout periods. Validate the 71-bus 2025 expansion with consistent 2025 infrastructure and operating inputs before validated DLR inference.\normalsize
\end{enumerate}
\par\medskip
\textcolor{teal}{\textbf{Current achievement:}} reproducible 2019 DC evidence and promising compact snapshots, with the dominant remaining uncertainties identified.
\source{Proposed next experiment, informed by [3,4]. No new shared-parameter fitting or independent holdout evaluation has been completed in this update.}
\end{frame}

\appendix
% 20
\begin{frame}{Appendix | Source reconstruction and time quality are audited}
\begin{columns}[T,onlytextwidth]
\column{0.48\textwidth}
\lead{Public-data reconstruction}
\begin{itemize}\tightitems
\item 12 NYISO monthly archives; 365 daily files; 1,897,236 raw rows.
\item 18 channels: seven interfaces and 11 schedules.
\item All 157,680 prepared hour/channel keys reproduced within $6.37\times10^{-12}$ in source units.
\end{itemize}
\column{0.48\textwidth}
\lead{8760 local-hour labels}
\begin{itemize}\tightitems
\item Two interpolated labels: March 10 at 02:00 and December 12 at 12:00.
\item November 3 at 01:00 combines both fall-DST folds.
\item Excluding these three labels changes WAPE from 11.2181\% to 11.2194\%.
\end{itemize}
\end{columns}
\par\bigskip
\textbf{Independent DC replay:} all 8,760 cases succeed; maximum branch-flow difference from MATPOWER is approximately $1.21\times10^{-11}$ MW.
\par\medskip
\minor{This numerical agreement verifies reproduction, not agreement with actual operation. Raw duplicates and irregular sampling were retained and audited.}
\source{[3,5]; NYISO source audit, channel summaries, and independent DC replay. Released data use timezone-naive local-hour labels.}
\end{frame}

% 21
\begin{frame}{Appendix | Five-hour matrix exposes input dependence}
\small
\begin{center}
\begin{tabular}{lrrrrrr}
\toprule
2019 local hour & A & B & C & D & C + Marcy & D + Marcy \\
\midrule
Jan 21, 18:00 & 9.40 & 16.96 & 19.52 & 19.73 & 14.43 & 16.97 \\
Apr 21, 04:00 & 22.62 & 36.66 & 27.91 & 39.50 & 19.08 & 29.85 \\
Jul 17, 16:00 & 20.72 & 26.23 & 30.83 & 25.40 & 26.10 & 27.53 \\
Jul 20, 16:00 & 12.34 & 16.13 & 23.05 & 19.42 & 16.84 & 17.43 \\
Jul 27, 02:00 & 18.86 & 26.87 & 22.84 & 22.94 & 23.20 & 27.95 \\
\bottomrule
\end{tabular}
\end{center}
\small
All entries are pooled WAPE (\%). A/B use the 46-bus reference; C/D use compact51. A/C use released paper inputs; B/D use current-pipeline 2019 inputs. Common paper-role cuts, DC reference 74.
\par\medskip
\begin{itemize}\tightitems
\item B maps the extra compact bus 9003 injection equally to 78 and 80 through an exact compact-network Kron injection map. This is a declared topology-dependent equivalent.
\item Released boundary targets use a raw-sample arithmetic mean; current preparation uses forward hold. Input differences include time aggregation.
\item July 17 at 08:00 is excluded from this matrix due to a current-input boundary gap; the released-input annual study retains the hour.
\end{itemize}
\source{[4]; network\_input\_matrix.csv, Kron injection map and validation, and input export description. Values rounded to two decimal places.}
\end{frame}

% 22
\begin{frame}{Appendix | Selecting by pooled error can hide one bad interface}
\begin{center}
\small% Generated from frozen study CSVs; requires booktabs.
\begin{tabular}{llcrr}
\toprule
Criterion & Rank & 2019 hour & Pooled & Worst \\
\midrule
pooled & 1 & Sep 06 02:00 & 1.77\% & 16.15\% \\
pooled & 2 & Sep 08 05:00 & 1.95\% & 5.59\% \\
pooled & 3 & Sep 06 03:00 & 2.17\% & 15.39\% \\
minimax & 1 & Jul 03 04:00 & 3.28\% & 4.85\% \\
minimax & 2 & Jul 22 04:00 & 2.63\% & 5.31\% \\
minimax & 3 & Sep 08 05:00 & 1.95\% & 5.59\% \\
\bottomrule
\end{tabular}

\end{center}
\par\medskip
\textbf{For preliminary examples, minimax selection is easier to defend:} every interface must perform reasonably well at the selected hour.
\par\medskip
\minor{The smallest pooled error is 1.77\%, but its worst interface error is 16.15\%. All rankings are retrospective; the low-error hours do not substitute for the 12.32\% annual compact + Marcy result.}
\source{[4]; best\_snapshots.csv. Pooled sums absolute mismatch / absolute observed flow; worst is the largest of seven individual proportional errors.}
\end{frame}

% 23
\begin{frame}{Appendix | Reproduction leaves physical limits to resolve}
\begin{itemize}\tightitems
\item \textbf{Released DC states are not bounded operating solutions.}\par
\small 8,745 hours exceed at least one positive RATE\_A value; 70 of 94 AC branches have no positive rating. DC convergence does not establish thermal or AC feasibility.\normalsize
\item \textbf{Source generation contains reconstruction assumptions.}\par
\small One July St. Lawrence prior is 921.12 MW against an 856 MW parameter capacity. The combined Zone-H generation prior assigned to nuclear proxies exceeds restored capacity in two of five development hours; it was retained transparently.\normalsize
\item \textbf{The inherited compact backbone is not identical to NYgrid.}\par
\small Original parent files contain 233 versus 227 branch records; 43--50 and 74--78 multiplicities differ, with additional S4 paths. No blanket branch pruning was applied.\normalsize
\item \textbf{The AC control checkpoint has a different scope.}\par
\small It restores two nuclear proxies with finite assumed reactive capability. Large reactive/voltage adjustments and April's loss increase require a better control prior before historical interpretation.\normalsize
\end{itemize}
\source{[3,4]; source capacity audit, branch-pair crosswalk, AC controls and fresh-PF validation. Prior compact S4/S7 calibration history is inherited; this experiment performs no new fit.}
\end{frame}

% 24
\begin{frame}{References and reproducibility}
\small
\begin{enumerate}\setlength{\itemsep}{0.75em}
\item M. V. Liu, B. Yuan, Z. Wang, J. A. Sward, K. M. Zhang, and C. L. Anderson. ``An Open Source Representation for the NYS Electric Grid to Support Power Grid and Market Transition Studies.'' \emph{IEEE Transactions on Power Systems}, 38(4), 3293--3303, 2023. \href{https://doi.org/10.1109/TPWRS.2022.3200887}{DOI: 10.1109/TPWRS.2022.3200887}.
\item \href{https://github.com/AndersonEnergyLab-Cornell/NYgrid/tree/47698b6c7823ae7b1bc6935e5601d5b64b8f918e}{Released NYgrid source}, pinned commit \texttt{47698b6}. This update reproduces its 2019 study with documented measurement corrections.
\item \href{https://github.com/yifeihsu/nyiso-dlr-testbed/blob/c472e250753b441223d8c6e029aa86e6dd56db8e/NYGRID_2019_INTERFACE_VERIFICATION.md}{Project: NYgrid 2019 interface verification}. Full-year metrics, source reconstruction, formula audit, and selected-hour records.
\item \href{https://github.com/yifeihsu/nyiso-dlr-testbed/blob/c472e250753b441223d8c6e029aa86e6dd56db8e/NYGRID_COMPACT_2019_COMPARISON.md}{Project: matched compact 2019 comparison}. Annual/matrix/ablation results, role crosswalks, source hashes, and separate AC checkpoint.
\item \href{https://mis.nyiso.com/public/P-32list.htm}{NYISO public P-32 archives}. Historical interface-flow and schedule records, audited against the released prepared inputs.
\end{enumerate}
\par\medskip
\minor{Scientific evidence is frozen at project commit \texttt{c472e25}. The accompanying LaTeX source, vector figures, editable CSVs, generation script, and hash manifests make this deck reviewable and rebuildable. No studies were rerun for this presentation.}
\end{frame}

\end{document}
